% Full added formal demonstrations for ED26-BM-RAT2021-10-1.
\begin{annexitem}{annex:p18.r101.aggregation}{thm:p18-r101-aggregation-failures}{Judgments, rankings, and plurality have different failure mechanisms}
\begin{proof}
On the agenda constrained by \(r\leftrightarrow(p\land q)\), the
three rows \((1,1,1),(1,0,0),(0,1,0)\) all satisfy the constraint.
Their majority values are \(p=1,q=1,r=0\), which falsifies it.
Thus coherent complete individual judgments do not suffice for
coherent issuewise-majority judgment. The premise-based rule accepts
the two majority premises and therefore accepts the conclusion; its
output \((1,1,1)\) is coherent but overrides two individual votes on
\(r\). This explicitly identifies the changed rule property.

For strict rankings \(a>b>c\), \(b>c>a\), and \(c>a>b\), the first
and third members rank \(a\) above \(b\), the first and second rank
\(b\) above \(c\), and the second and third rank \(c\) above \(a\).
The majority relation therefore contains a strict directed cycle. Each
input ranking is transitive; the output is not.

Before splitting, sixty voters choose \(a\) and forty choose \(b\),
so plurality selects \(a\). Replace \(a\) by variants
\(a_1,a_2,a_3\), with each receiving twenty of its former supporters,
and retain the forty votes for \(b\). Then \(b\) is the unique
plurality winner because forty exceeds twenty. If each of those sixty
voters instead approves all three variants and the forty remaining
voters approve only \(b\), all variants receive sixty approvals and
\(b\) receives forty. This establishes avoidance of this particular
splitting example, not a general dominance result for approval.
\end{proof}
\end{annexitem}

\begin{annexitem}{annex:p18.r101.pooling}{thm:p18-r101-pooling}{Finite pools, common learning, and refinement}
\begin{proof}
Nonnegativity and \(L(\Omega)=\sum_iw_i=1\) follow immediately.
For disjoint \(A,B\), finite summation gives
\(L(A\cup B)=\sum_iw_i(P_i(A)+P_i(B))=L(A)+L(B)\).
The same identity shows that summing over the worlds in any coarse cell
before or after linear pooling agrees.

For the displayed numerical example, the linear pool is
\((0.4,0.3,0.3)\), so conditioning on \(E=\{a,b\}\) assigns
\(0.4/0.7=4/7\) to \(a\). The separate conditionals assign
\(0.6/0.9=2/3\) and \(0.2/0.5=2/5\), whose equal-weight average
is \(8/15\). These are unequal because \(4\cdot15=60\) and
\(8\cdot7=56\). All conditioning events have positive probability.

For the geometric pool, its denominator is positive and finite, so its
mass function is well defined and normalized. For \(\omega\in E\),
pooling the member conditionals gives unnormalized weight
\[
 \prod_i\left(\frac{P_i(\omega)}{P_i(E)}\right)^{w_i}
 =\frac{\prod_iP_i(\omega)^{w_i}}{\prod_iP_i(E)^{w_i}}.
\]
The denominator on the right is constant over \(\omega\in E\) and
cancels when the weights on \(E\) are normalized. The resulting mass
is therefore the original \(G\) conditioned on \(E\). Worlds outside
\(E\) receive zero after conditioning and need not enter the product
formula for its positive support. This proves the assertion for fixed
weights and a common learned event, not arbitrary private learning.

For refinement, both coarse inputs are \((1/2,1/2)\), hence their
geometric pool is the same. After splitting \(A\), the two refined
inputs are \((0.45,0.05,0.5)\) and \((0.05,0.45,0.5)\).
Their equal-weight geometric weights are
\((\sqrt{0.45\cdot0.05},\sqrt{0.05\cdot0.45},0.5)
=(0.15,0.15,0.5)\), with total \(0.8\). Consequently the aggregate
mass of \(A\) is \(0.3/0.8=3/8\), not \(1/2\). Positivity holds
on both spaces, so the discrepancy is not an undefined zero-support case.
\end{proof}
\end{annexitem}

\begin{annexitem}{annex:p18.r101.jury}{thm:p18-r101-jury}{Classical binomial majority and finite conditional mixtures}
\begin{proof}
For \(n=2m+1\) independent Bernoulli indicators with success
probability \(t\), exactly \(k\) successes has probability
\({n\choose k}t^k(1-t)^{n-k}\). Summing over \(k\geq m+1\)
proves the formula.

Couple a group of \(2m+1\) with two new independent voters. An old
incorrect majority becomes correct only when the old number of successes
is \(m\) and both new voters succeed. An old correct majority becomes
incorrect only when the old number is \(m+1\) and both new voters fail.
For all other old success counts the majority status cannot change.
Thus, using \({2m+1\choose m}={2m+1\choose m+1}\),
\[
 M_{2m+3}(t)-M_{2m+1}(t)
 ={2m+1\choose m}t^{m+1}(1-t)^{m+1}(2t-1).
\]
This is positive for \(1/2<t<1\), negative for \(0<t<1/2\),
and zero at \(0,1/2,1\).

For \(t>1/2\), write \(\delta=t-1/2>0\). Independence gives
\(\mathbb E(S_n/n)=t\) and
\(\mathbb E[(S_n/n-t)^2]=t(1-t)/n\). On the event
\(|S_n/n-t|\geq\delta\), the squared deviation is at least
\(\delta^2\). Nonnegativity therefore yields
\[
 \Pr(|S_n/n-t|\geq\delta)
 \leq \frac{t(1-t)}{n\delta^2}.
\]
Incorrect majority is contained in this event, so its probability tends
to zero. This derives the required bound rather than assuming a law of
large numbers. For \(t<1/2\) the same reasoning bounds the probability
of a correct majority and makes it tend to zero. At \(t=1/2\),
complementing all votes bijects correct and incorrect outcomes with equal
probabilities; odd size excludes a tie. Thus \(M_n(1/2)=1/2\).
At \(t=0,1\) the limits are immediate.

Condition on each finite problem value and apply the binomial result.
The law of total probability gives
\(\sum_j\mu_jM_n(t_j)\); passing a limit through this finite sum
gives the displayed limit. Combine problem values with the same
competence into their total probability mass. Pair each value
\(t=1/2+d\) with \(1-t\), assigning zero mass to an absent partner.
The increment formula at the two values has equal magnitude and
opposite sign. If the upper mass is at least the lower mass, the
pair's contribution to the mixture increment is nonnegative. The
central and endpoint contributions are zero. Summing the finitely many
pairs proves weak increase. Strict increase needs positive excess mass
at some upper competence strictly between one half and one; endpoint
excess alone does not suffice.

In the fallible-improving example, competence values \(0.8,0.2\)
have masses \(0.8,0.2\). Their reflected dominance holds, individual
accuracy is \(0.8^2+0.2^2=0.68\), and the majority limit is \(0.8\).
In the mean-only counterexample, competence values \(0.99,0.49\)
each have mass \(1/2\), so individual accuracy is \(0.74\), while
the limit is \(1/2\). There is no mass at \(0.51\) to dominate
the lower-half mass at \(0.49\); hence the distributional premise
fails. These calculations establish the stated boundary and do not
identify the parameters of any observed electorate.
\end{proof}
\end{annexitem}
